\documentclass[12pt]{article}
\usepackage{fontspec}
\usepackage[a4paper,top=2.3cm,bottom=2.3cm,left=2.5cm,right=2.5cm]{geometry}
\usepackage{xcolor}
\IfFontExistsTF{PT Serif}{\setmainfont{PT Serif}}{%
  \IfFontExistsTF{Times New Roman}{\setmainfont{Times New Roman}}{\setmainfont{DejaVu Serif}}}
\IfFontExistsTF{Menlo}{\setmonofont{Menlo}}{}
\usepackage[english,russian]{babel}
\definecolor{barcol}{HTML}{3B6EA5}
\definecolor{rulecol}{gray}{0.55}
\definecolor{headcol}{HTML}{1F3A5F}
\definecolor{spokencol}{HTML}{1A2A3A}
% Spoken-text block: indented + colored (package-free; page-breakable).
% A left rule is drawn beside the first line of each paragraph via \everypar.
\newenvironment{spoken}{%
  \par\addvspace{0.45em}\begingroup
  \color{spokencol}%
  \setlength{\leftskip}{1.5em}\setlength{\rightskip}{0.2em}%
  \setlength{\parindent}{0pt}%
  \everypar{\llap{{\color{barcol}\vrule width 2pt height 0.9em depth 0.35em}\hspace{0.9em}}}%
}{%
  \par\endgroup\addvspace{0.45em}}
\setlength{\parindent}{0pt}
\setlength{\parskip}{0.5em}
\setcounter{secnumdepth}{0}
\emergencystretch=3em
\sloppy
\frenchspacing
\begin{document}

\begin{center}{\LARGE\bfseries\color{headcol} Текст устного доклада — защита ВКР}\end{center}\vspace{0.4em}

\begin{center}\textbf{Можаев Роман Михайлович} · магистратура ФПМИ МФТИ, базовая кафедра ИППИ РАН \\ \textbf{Тема:} «Исследование глобальных квазиньютоновских методов оптимизации для задач оптимизации и вариационных неравенств» \\ \textbf{Научный руководитель:} Д. И. Камзолов, к.ф.-м.н. · \textbf{Рецензент:} А. В. Рогозин, к.ф.-м.н. \\ \textbf{Защита:} 16.06.2026\end{center}

\medskip\noindent{\color{rulecol}\rule{\linewidth}{0.4pt}}\par\medskip

\subsection*{\color{headcol} Как пользоваться этим файлом}

\begin{itemize}\setlength{\itemsep}{0.2em}
\item \textbf{Тайминг доклада:} рассчитан на \textbf{\textasciitilde{}14 минут}. Под каждым слайдом — целевое время; сумма основного хода \ensuremath{\approx} 14 мин.
\item \textbf{Доказательство теоремы 2.10} вынесено из колоды; при вопросе изложить устно (см. блок вопросов — «Откуда взялось равенство…»).
\item Если регламент \textbf{окажется 10–12 минут} — сжать слайд 10 (Следствие, помечен {\ttfamily\small\color{rulecol}[можно сжать]}) и слайд 13 (Discrete BVP) до одной фразы.
\item Это \textbf{разговорный} текст, а не чтение слайда вслух. Цифры и формулы — на экране; голосом даём «зачем» и связки между слайдами.
\item После доклада идут \textbf{ответы на замечания} и \textbf{заготовки на вопросы} — отдельные блоки, читаются по ситуации.
\end{itemize}

\medskip\noindent{\color{rulecol}\rule{\linewidth}{0.4pt}}\par\medskip

\section*{\color{headcol} ЧАСТЬ 1. ДОКЛАД (по слайдам)}

\subsubsection*{\color{headcol} Титул  ·  \textasciitilde{}25 сек}
\begin{spoken}
Уважаемые члены комиссии, добрый день. Меня зовут Роман Можаев. Тема моей выпускной квалификационной работы — исследование квазиньютоновских методов с памятью для систем нелинейных уравнений, задач оптимизации и вариационных неравенств. Научный руководитель — Дмитрий Игоревич Камзолов. Позвольте перейти к докладу.\par
\end{spoken}

\subsubsection*{\color{headcol} Слайд 1 — План доклада  ·  \textasciitilde{}20 сек}
\begin{spoken}
Доклад построю так: сначала мотивация, цель и предмет работы; затем основной теоретический результат — точная теорема об ограниченной деградации для мультисекущего Бройдена; далее его limited-memory вариант; симметричный аналог для оптимизации; и в конце — приложение к вариационным неравенствам. Завершу выводами и перспективами.\par
\end{spoken}

\subsubsection*{\color{headcol} Слайд 2 — Мотивация  ·  \textasciitilde{}50 сек}
\begin{spoken}
Базовая задача вычислительной математики — решить нелинейную систему \(F(x)=0\), либо найти минимум гладкой функции. Такие задачи возникают буквально всюду: дискретизация дифференциальных уравнений, обучение моделей, поиск равновесий в играх и экономике.\par
Замкнутого решения в общем случае нет — строится последовательность приближений \(x_{0}, x_{1}, x_{2}, \dots \). Когда размерность \(n\) велика, \textbf{решающей становится стоимость одной итерации}: считать и обращать производную на каждом шаге слишком дорого. Именно здесь работают квазиньютоновские методы — к ним и перехожу.\par
\end{spoken}

\subsubsection*{\color{headcol} Слайд 3 — От Ньютона к КН методам (1/2): идея метода  ·  \textasciitilde{}30 сек}
\begin{spoken}
Эталон точности — метод Ньютона: он строит линейную модель оператора в текущей точке и делает шаг. Сходится быстро, но на каждой итерации требует \textbf{якобиан} и решение линейной системы — это дорого, а часто якобиан вообще недоступен.\par
Квазиньютоновская идея: заменить настоящий якобиан матрицей \(B_{k}\), которую мы \textbf{достраиваем по уже наблюдённым данным}, без вычисления производных. Платим за итерацию дёшево, а скорость стараемся удержать близкой к ньютоновской.\par
\end{spoken}

\subsubsection*{\color{headcol} Слайд 4 — От Ньютона к КН методам (2/2): проблема и BD  ·  \textasciitilde{}28 сек}
\begin{spoken}
И вот ключевая проблема, ради которой всё затевается \textit{(показать на «Проблему»)}: классические обновления — Бройден, SR1, BFGS — сходятся \textbf{лишь локально} и чувствительны к старту; без глобализации сходимость не гарантирована. Это классические результаты — Бройден–Деннис–Море, семьдесят третий год, и Гриванк, восемьдесят шестой. Локальную сходимость выводят из \textit{ограниченной деградации} — оценки того, насколько ошибка \(E_{k} = B_{k} - J^{*}\) может вырасти за шаг. Это \textbf{одна из двух} классических компонент локальной теории (вторая — условие Денниса–Море для сверхлинейной скорости), а не «устойчивость» сама по себе и не глобальная гарантия. Именно эту оценку работа и уточняет — дальше покажу, как окно прошлых шагов контролирует деградацию.\par
\end{spoken}

\subsubsection*{\color{headcol} Слайд 5 — Проблема численно: чувствительность к старту  ·  \textasciitilde{}35 сек  ·  {\ttfamily\small\color{rulecol}[можно сжать/пропустить — факт процитирован на сл. 4]}}
\begin{spoken}
Чтобы это не звучало голословно — численная иллюстрация, и сознательно на \textbf{стандартной} задаче из набора Море–Гарбоу–Хиллстрома, а не на своём примере. Классический Бройден, причём с \textbf{точным} начальным якобианом — то есть в самых выгодных для метода условиях. Зелёные траектории — старт близко к решению: сверхлинейная сходимость за десяток итераций. Красные — тот же метод, старт подальше: \textbf{расходится}. Вывод ровно теоретический — область сходимости конечна, и без глобализации сходимость не гарантирована. Именно эту деградацию дальше и берётся контролировать окно прошлых шагов.\par
\end{spoken}

\subsubsection*{\color{headcol} Слайд 6 — Цель и предмет  ·  \textasciitilde{}50 сек  \textit{(важно проговорить «глобальные»)}}
\begin{spoken}
Цель работы — тройная: дать \textbf{количественный} анализ мультисекущего Бройдена для систем \(F(x)=0\); построить его \textbf{симметричный} аналог для гладкой оптимизации; и показать роль точности якобиана во внешнем методе для вариационных неравенств. Предмет — квазиньютоновские обновления, которые держат не одно, а \textbf{\(p\) прошлых секущих условий} сразу, то есть помнят несколько последних шагов.\par
Сразу сниму возможный вопрос по названию \textit{(устно)}. Слово \textbf{«глобальные»} здесь — в двух смыслах: память о прошлых секущих, расширенное окно; и внешний глобальный метод второго порядка для вариационных неравенств в четвёртой главе. \textbf{Речь не о теореме глобальной сходимости} — работа сознательно локальная по теории, и это я ещё подчеркну.\par
А как именно окно прошлых шагов контролирует деградацию — на следующем слайде.\par
\end{spoken}

\subsubsection*{\color{headcol} Слайд 7 — Мультисекущая идея  ·  \textasciitilde{}60 сек}
\begin{spoken}
Теперь содержательно. Секущее условие — это требование, чтобы новая матрица правильно действовала вдоль сделанного шага: \(B_{k+1} s_{k} = y_{k}\), где \(s\) — приращение точки, \(y\) — приращение оператора. По сути это дискретный аналог производной вдоль шага.\par
Классический Бройден держит \textbf{ровно одно} такое условие и при этом меняет матрицу минимально — принцип наименьшего изменения. Геометрически это проектор ранга один: гарантия только вдоль одного направления.\par
Моя идея \textit{(показать на нижнюю формулу)}: держать \textbf{окно} из нескольких последних шагов \(S_{p}\) и потребовать секущее условие сразу для всех столбцов: \(B S_{p} = Y_{p}\). Решение того же вариационного принципа — вот эта явная формула. Теперь проектор ограниченной деградации расширяется до \textbf{ранга \(p+1\)} на всём подпространстве шагов. При \(p=0\) всё схлопывается обратно в классический Бройден — то есть это честное обобщение, а не другой метод.\par
\end{spoken}

\subsubsection*{\color{headcol} Слайд 8 — Теорема (гл.2) (1/2): равенство и bounded deterioration  ·  \textasciitilde{}50 сек  \textit{(главный результат, не торопиться)}}
\begin{spoken}
Это центральный результат работы. \(E_{k} = B_{k} - J^{*}\) — ошибка аппроксимации якобиана; \(\Pi _{k}\) — ортопроектор на подпространство накопленных шагов.\par
Пункт \textbf{(a)} — \textbf{точное равенство}: норма ошибки на следующем шаге равна норме на текущем, \textbf{минус} проекция ошибки на окно, \textbf{плюс} малый остаточный член \(\Delta _{k}\). Подчеркну: это \textbf{равенство, а не оценка сверху} — мы видим точный баланс ошибки, без запаса. Аналитически это сильнее: вместо «не хуже, чем» мы знаем «ровно столько».\par
\end{spoken}

\subsubsection*{\color{headcol} Слайд 9 — Теорема (гл.2) (2/2): константа $C_{PB}$, не зависящая от $n$  ·  \textasciitilde{}45 сек}
\begin{spoken}
Пункт \textbf{(b)} — отсюда ограниченная деградация: остаточный член контролируется как \(C_{\mathrm{PB}}^{2}\cdot \rho _{k}^{2}\), где \(\rho _{k}\) — расстояние до корня. И ключевое — вот эта константа \(C_{\mathrm{PB}}\): в неё входят липшицева константа, корень из \(p+1\) и \(\kappa _{p}\) — \textbf{обусловленность окна шагов}, то есть насколько направления в окне линейно независимы. Самое важное: \textbf{константа не зависит от размерности \(n\)}. Именно это открывает дорогу к limited-memory варианту для огромных задач.\par
И ещё одно наблюдение \textit{(нижний пункт)}: отрицательный член — тот, что «съедает» ошибку, — у мультисекущей версии \textbf{не меньше}, чем у классического Бройдена. То есть окно может только помочь, не навредив.\par
\end{spoken}

\subsubsection*{\color{headcol} Слайд 10 — Следствие: контроль ошибки на подпространстве  ·  \textasciitilde{}50 сек  ·  {\ttfamily\small\color{rulecol}[можно сжать]}}
\begin{spoken}
Прямое следствие центральной теоремы. Матрица \(B_{k+1}\) приближает настоящий якобиан с точностью \(O(\rho _{k})\) \textbf{на всём подпространстве} прошлых шагов, причём эта точность зависит только от липшицевой константы и расстояния до корня — \textbf{но не от текущей величины ошибки} \(\|E_{k}\|\). Это и есть структурное преимущество: у классического Бройдена такой безусловный контроль есть только вдоль \textbf{одного} последнего направления, а на остальном подпространстве ошибка определяется всей предысторией итераций. Окно превращает одномерную гарантию в многомерную — и дальше это окупится в главе про вариационные неравенства.\par
\end{spoken}

\subsubsection*{\color{headcol} Слайд 11 — L-PB (1/2): почему обновляем обратную ранг-1  ·  \textasciitilde{}35 сек}
\begin{spoken}
Зачем нужна была независимость от \(n\). Плотная матрица \(B_{k}\) — это \(O(n^{2})\) памяти и \(O(n^{3})\) операций на шаг. При \(n\) порядка ста тысяч это около восьмидесяти гигабайт — метод просто неприменим.\par
Выход: обновлять можно сразу \textbf{обратную} матрицу, и это обновление — \textbf{ранг один}, по формуле Шермана–Моррисона. Весь фокус в специальном направлении \(v_{k}\): оно единично вдоль текущего шага и \textbf{ортогонально всем прошлым} шагам окна — поэтому ранг-один поправка \textbf{добавляет} новое секущее условие, \textbf{не разрушая} уже накопленные. Именно из-за этого \(v_{k}\) мы храним не пары, как L-BFGS, а \textbf{тройки} \((s, y, v)\) — скользящее окно из \(m\) штук; само применение обратной матрицы идёт рекурсией, как в L-BFGS. Стоимость — \(O(nm)\) памяти и \(O(nm^{2})\) операций, линейно по \(n\). При \(p_{\mathrm{max}} = 0\) это в точности классический limited-memory Бройден.\par
\end{spoken}

\subsubsection*{\color{headcol} Слайд 12 — L-PB (2/2): limited memory и перенос теоремы  ·  \textasciitilde{}30 сек}
\begin{spoken}
Глубина окна растёт до предела \(p_{\mathrm{max}}\), и троек нужно хранить хотя бы \(p_{\mathrm{max}} + 1\). Тогда — это Утверждение 2.1 — limited-memory форма \textbf{совпадает} с плотным проектированным Бройденом на последних \(p_{\mathrm{max}}+1\) секущих, а значит центральная теорема переносится на L-PB \textbf{дословно}, с той же константой \(C_{\mathrm{PB}}\). Она зависит от обусловленности окна, а не от размерности — то есть экономия памяти \textbf{ничего не стоит в теории}: гарантия сохраняется точно, а не приближённо.\par
\end{spoken}

\subsubsection*{\color{headcol} Слайд 13 — Численный результат: Discrete BVP (n=100)  ·  \textasciitilde{}50 сек  \textit{(ключевой месседж)}}
\begin{spoken}
Теперь подтверждение на практике. Задача — дискретная краевая задача, размерность сто, двадцать случайных стартов. Красная линия — классический Бройден: он \textbf{не сходится ни на одном} из двадцати стартов. Стоит включить окно хотя бы глубины один — метод \textbf{устойчиво сходится на всех}.\par
Обратите внимание на тонкость, которая прямо иллюстрирует теорию: ошибка у классического Бройдена остаётся \textbf{ограниченной} — теория ограниченной деградации не нарушается, — но сходимости при этом нет. То есть «ошибка ограничена» и «метод сходится» — это разные вещи. Окно переводит метод из режима «ошибка ограничена, но стоим на месте» в режим \textbf{сверхлинейной сходимости}.\par
\end{spoken}

\subsubsection*{\color{headcol} Слайд 14 — L-PB в высокой размерности: Broyden Banded  ·  \textasciitilde{}45 сек}
\begin{spoken}
Проверка масштабируемости. Размерность — \textbf{сто тысяч}, где плотная форма потребовала бы те самые восемьдесят гигабайт и физически невозможна. L-PB сходится, и — это главное на графике — число итераций \textbf{не растёт с ростом размерности}. Деградация по глубине окна \(m\) умеренная и согласуется с теорией, ровно как ведёт себя L-BFGS в задачах оптимизации. То есть метод действительно переносится в большую размерность, а не только формально.\par
\end{spoken}

\subsubsection*{\color{headcol} Слайд 15 — Симметричный случай (1/2): зачем симметрия и в чём трудность  ·  \textasciitilde{}30 сек}
\begin{spoken}
Переходим к оптимизации. Здесь приближаем не якобиан, а \textbf{гессиан} — а он симметричен, и естественно хотеть, чтобы и \(B_{k}\) была симметричной: тогда матрица отвечает корректной квадратичной модели. Но есть препятствие — лемма Шнабеля: в общем положении \textbf{нельзя одновременно} сохранить симметрию и несколько секущих условий. Они конфликтуют. Идея — взять базовое симметричное обновление, которое точно держит \textbf{текущую} секущую, и потратить единственную оставшуюся степень свободы на согласование с \textbf{прошлыми}.\par
\end{spoken}

\subsubsection*{\color{headcol} Слайд 16 — Симметричный случай (2/2): конструкция SS-U  ·  \textasciitilde{}30 сек}
\begin{spoken}
Конкретно: берём базовое симметричное обновление — SR1 или PSB; оно уже точно держит текущую секущую, то есть \(B_{k}' s_{k} = y_{k}\). Смотрим на его невязку на окне: для прошлых шагов \(B_{k}' s_{j} - y_{j}\) в общем не ноль, а для текущего столбца — ровно ноль. Добавляем \textbf{ранг-1 коррекцию в направлении, ортогональном} \(s_{k}\). Поскольку коррекция ортогональна \(s_{k}\), она \textbf{не портит текущее секущее условие} и \textbf{сохраняет симметрию}, а её единственную степень свободы мы тратим на то, чтобы максимально погасить невязку по \textbf{прошлым} секущим. По теореме 3.8 это даёт сразу: точное текущее секущее условие, симметрию, ранг коррекции один и монотонное убывание накопленной невязки по прошлым секущим.\par
\end{spoken}

\subsubsection*{\color{headcol} Слайд 17 — Сходимость SS-U (1/2): ограниченная деградация  ·  \textasciitilde{}28 сек}
\begin{spoken}
Что со сходимостью. Лемма 3.12 — композитная ограниченная деградация: норма ошибки убывает на проекцию вдоль шага плюс остаточный член \(\beta _{k}\). Что такое \(\beta _{k}\) — это суммарная деградация за шаг: вклад базового обновления и ранг-1 коррекции в рост ошибки; сумма этих остатков по всем шагам конечна. И отсюда итерация сходится \textbf{Q-линейно} — подчеркну, это \textbf{выходное} свойство: мы его не предполагаем, а получаем из самого неравенства при хорошем старте.\par
\end{spoken}

\subsubsection*{\color{headcol} Слайд 18 — Сходимость SS-U (2/2): от деградации к сверхлинейности  ·  \textasciitilde{}28 сек  \textit{(честно про «условность»)}}
\begin{spoken}
Зачем нужна связка между леммой и теоремой. Просуммируем неравенство по всем шагам: слева норма ошибки неотрицательна, остатки \(\beta _{k}\) суммируемы — значит, сумма проекций ошибки вдоль шагов конечна, а тогда сами проекции стремятся к нулю. Но это ровно \textbf{условие Денниса–Море} — классическая посылка сверхлинейной сходимости. То есть связка показывает, что \textbf{посылка теоремы 3.9 выполняется автоматически}, без отдельного предположения. А сама теорема 3.9 — \textbf{условная} Q-сверхлинейность: при невырожденном пределе и допустимости полного шага из условия Денниса–Море следует сверхлинейная скорость. Я специально подчёркиваю «условная» — это локальный результат при явно названных предположениях, и я их называю честно, не выдавая за безусловную теорему. Где границы — видно и в эксперименте, и в приложении про глобализацию.\par
\end{spoken}

\subsubsection*{\color{headcol} Слайд 19 — Численный результат: SS-PSB спасает PSB  ·  \textasciitilde{}50 сек  \textit{(честная граница)}}
\begin{spoken}
Эксперимент к третьей главе — и здесь я хочу быть предельно честным насчёт границ применимости. Чистый PSB деградирует: на расширенном Розенброке сходится примерно в двадцати из пятидесяти случаев. SS-PSB — устойчиво, сорок восемь–пятьдесят из пятидесяти.\par
\textbf{Но} — и это написано прямо на слайде — SS-коррекция помогает там, где базовое обновление \textbf{реально деградирует}. Для SR1 со skip-rule деградации нет — и там SS-SR1 выигрыша \textbf{не даёт}, я это тоже показываю в приложении. Это не недостаток метода, а \textbf{осознанно очерченная область применимости}: я показываю и где метод выигрывает, и где он нейтрален.\par
\end{spoken}

\subsubsection*{\color{headcol} Слайд 20 — Приложение к вариационным неравенствам  ·  \textasciitilde{}65 сек  \textit{(честно про соавторство)}}
\begin{spoken}
Четвёртая глава — приложение. Вариационные неравенства — это общий язык для седловых задач, равновесий и, например, обучения с противником; для них тоже строят методы второго порядка.\par
Здесь рассматривается метод \textbf{VIJI} для монотонных вариационных неравенств; он использует линейную модель оператора с \textbf{неточным} якобианом, и вся неточность сворачивается в одно число \(\delta \) — норму разности настоящего якобиана и его приближения.\par
Сразу обозначу авторство честно: \textbf{оценка скорости} — вот эта формула — получена в \textbf{соавторской работе}, принятой на NeurIPS 2024; зависимость от \(\delta \) там неулучшаема, есть нижняя оценка.\par
А \textbf{мой собственный вклад в эту главу} — Утверждение 4.1: я \textbf{подставляю} в роль \(\delta \) мою мультисекущую оценку из второй главы. Получается \(\delta _{k}  \le  \|E_{k}\| + L_{1}\cdot \rho _{k}\) — с \textbf{подпространственным} \(O(\rho _{k})\)-контролем из следствия, вместо грубого равномерного потолка \(\delta   \le  2L_{0}\), который даёт обычный односекущий L-Broyden. То есть локальная теория глав 2–3 напрямую улучшает гарантию внешнего метода — локально, вблизи решения (глобального переноса я не утверждаю). Это и есть смысловой мост всей работы.\par
\end{spoken}

\subsubsection*{\color{headcol} Слайд 21 — Заключение  ·  \textasciitilde{}55 сек}
\begin{spoken}
Подытожу. Получены: точное разложение ограниченной деградации и \textbf{не зависящая от размерности} константа для мультисекущего Бройдена; количественный \(O(\rho _{k})\)-контроль ошибки на подпространстве; его limited-memory вариант L-PB, работающий при \(n\) порядка ста тысяч; симметричное обновление SS-U со сверхлинейной сходимостью, эмпирически стабилизирующее деградирующие SR1 и PSB; и подстановка локальной оценки в метод для вариационных неравенств.\par
Перспективы — формальная глобализация через line search, крупномасштабные задачи вроде PINN и равновесных моделей, и систематическое сравнение с L-BFGS и Anderson. Спасибо.\par
\end{spoken}

\subsubsection*{\color{headcol} Спасибо  ·  \textasciitilde{}10 сек}
\begin{spoken}
Спасибо за внимание. Готов ответить на вопросы.\par
\end{spoken}

\textbf{Итого \ensuremath{\approx} 14 минут.} Чтобы уйти в 12 — сжать слайд 10 (Следствие) и слайд 13 (Discrete BVP) до одной фразы.

\medskip\noindent{\color{rulecol}\rule{\linewidth}{0.4pt}}\par\medskip

\section*{\color{headcol} ЧАСТЬ 2. ОТВЕТЫ НА ЗАМЕЧАНИЯ РЕЦЕНЗЕНТА (А. В. Рогозин)}

\begin{spoken}
Общая интонация: \textbf{поблагодарить, согласиться по существу, показать, что замечание либо уже учтено в докладе, либо осознанно вынесено в перспективы.} Все замечания рецензента — рекомендательные, оценка «отлично».\par
Отдельного backup-слайда с ответами рецензенту в колоде нет — отвечать устно по пунктам ниже.\par
\end{spoken}

\textbf{Вступительная фраза:}
\begin{spoken}
Благодарю рецензента за внимательный разбор и конструктивные замечания. Позвольте ответить по пунктам.\par
\end{spoken}

\textbf{Замечание 1 — терминология во введении («геометрия окна», «обусловленность окна шагов», «таксономия Fang–Saad»).}
\begin{spoken}
Согласен, это узкоспециальные термины. Поясню их: «обусловленность окна» \(\kappa _{p}\) — это число обусловленности матрицы накопленных шагов \(S_{p}\); оно измеряет, насколько направления в окне линейно независимы, и именно оно входит в мою константу. «Таксономия Fang–Saad» — это классификация мультисекущих обновлений по тому, какое условие и в какой метрике минимизируется. В докладе я эти термины либо ввожу явно, либо не использую. Для будущей журнальной публикации замечание учту — добавлю пояснения и постановочный глоссарий во введение.\par
\end{spoken}

\textbf{Замечание 2 — во введении нужна точная постановка задачи и идеи метода в виде формул.}
\begin{spoken}
Принимаю. Отмечу, что в докладе это сделано: постановка \(F(x)=0\) и \(min f(x)\) — на слайде мотивации, а ключевая идея — мультисекущее наименьшее изменение и расширенный проектор — выписана формулами на слайде «Мультисекущая идея». Это же усиление я перенесу во введение печатного текста при подготовке к публикации.\par
\end{spoken}

\textbf{Замечание 3 — эксперименты усилить сравнением с более разнообразным набором современных алгоритмов.}
\begin{spoken}
Справедливо. Поясню логику текущих экспериментов: они построены так, чтобы \textbf{изолировать эффект окна} внутри семейства Бройдена — классический Бройден против PB и L-PB, — и честно очертить границу SS-коррекции. Сравнение с L-BFGS уже обозначено на слайде высокой размерности, Anderson — в обзоре литературы работы. Систематический бенчмарк против L-BFGS, Anderson и нелинейного CG — это следующий естественный шаг, он в перспективах; но он не входит в \textbf{защищаемые теоретические положения}, которые носят аналитический характер.\par
\end{spoken}

\textbf{Замечание 4 — сверхлинейность SS-U при сильных локальных предположениях ограничивает применимость.}
\begin{spoken}
Полностью согласен, и я это специально подчёркиваю как в тексте, так и в докладе. Теорема 3.9 — \textbf{условная} Q-сверхлинейность: при невырожденном пределе и допустимости полного шага. Это сознательно \textbf{локальный} результат. Более того, в приложении я показываю, что без глобализации полный квазиньютоновский шаг расходится даже вблизи решения на невыпуклых задачах — то есть границы применимости очерчены явно и честно, а не замаскированы. Снятие этих предположений — через line search — это и есть направление формальной глобализации в перспективах.\par
\end{spoken}

\textbf{Замечание 5 — практическая эффективность в большой размерности требует более масштабных экспериментов.}
\begin{spoken}
Согласен. На сегодня показан модельный эксперимент при \(n = 10^{5}\) — там, где плотная форма физически невозможна, а L-PB сходится за число итераций, не растущее с размерностью. Это \textbf{доказательство принципа}. Индустриальные крупномасштабные бенчмарки — PINN, равновесные модели — действительно требуют отдельной большой работы и вынесены в перспективы. Подчеркну, что основной вклад работы — теоретический: именно независимость константы от \(n\) создаёт \textbf{обоснование} для таких экспериментов.\par
\end{spoken}

\textbf{Закрывающая фраза:}
\begin{spoken}
Ещё раз благодарю рецензента; замечания принимаю и считаю их точным указанием направления дальнейшей работы.\par
\end{spoken}

\medskip\noindent{\color{rulecol}\rule{\linewidth}{0.4pt}}\par\medskip

\section*{\color{headcol} ЧАСТЬ 3. ОТВЕТЫ НА ЗАМЕЧАНИЯ НАУЧНОГО РУКОВОДИТЕЛЯ (Д. И. Камзолов)}

\begin{spoken}
Руководитель сам пишет, что слабости «носят рекомендательный характер и не снижают оценки». Отвечать развёрнуто не нужно — достаточно короткого согласия, \textbf{если} прозвучит вопрос.\par
\end{spoken}

\textbf{Замечание 1 — отсутствие теорем о глобальной сходимости (нужен более сложный внешний метод VIQA).}
\begin{spoken}
Согласен. Работа по теории сознательно локальная. Слово «глобальные» в названии относится к памяти прошлых секущих и к внешнему глобальному методу для вариационных неравенств, а не к теореме глобальной сходимости — это я оговариваю на слайде «Цель». Формальная глобализация через line search — первый пункт перспектив.\par
\end{spoken}

\textbf{Замечание 2 — расширить эксперименты новыми методами и задачами.}
\begin{spoken}
См. ответ на замечание 3 рецензента — то же направление; принимаю.\par
\end{spoken}

\medskip\noindent{\color{rulecol}\rule{\linewidth}{0.4pt}}\par\medskip

\section*{\color{headcol} ЧАСТЬ 4. ВЕРОЯТНЫЕ ВОПРОСЫ КОМИССИИ И ЗАГОТОВКИ ОТВЕТОВ}

\textbf{В: Почему всё-таки в названии «глобальные», если глобальной сходимости нет?}
\begin{spoken}
Два смысла, оба корректны: глобальная \textbf{память} (окно прошлых секущих, в отличие от локального одношагового Бройдена) и внешний \textbf{глобальный метод второго порядка} для вариационных неравенств в главе 4. Теоремы глобальной сходимости в смысле «из любой точки» я не утверждаю и это явно оговариваю.\par
\end{spoken}

\textbf{В: Чем это принципиально отличается от Block Broyden (NeurIPS 2023) или Anderson acceleration?}
\begin{spoken}
Block Broyden использует \textbf{случайный} блок направлений и даёт ухудшение «в среднем», без константы обусловленности. У меня — \textbf{детерминированное окно вдоль траектории} и \textbf{точное равенство} с явной \(\kappa _{p}\). Anderson — это Type-II мультисекущий метод, но без явной модели якобиана и без покомпонентной оценки ограниченной деградации, которую я доказываю. Различия готов детализировать устно — они разобраны в обзоре литературы работы.\par
\end{spoken}

\textbf{В: Откуда взялось «равенство, а не неравенство», и что это даёт?}
\begin{spoken}
Из ортогонального построчного разложения ошибки и теоремы Пифагора: строки \(E_{k} \Pi \perp \) лежат в ортогональном дополнении к окну, строки остатка \(R_{k}\) — в самом окне, они ортогональны. Равенство даёт \textbf{точный баланс}: видно, сколько ошибки снимается проекцией и сколько добавляется остатком, без запаса. У Гэя–Шнабеля был только верхний барьер. \textit{(Доказательство из колоды убрано; готов изложить эти пять строк устно или по тексту работы.)}\par
\end{spoken}

\textbf{В: Что если окно плохо обусловлено — \(\kappa _{p}\) большое?}
\begin{spoken}
Тогда константа растёт — и это честно отражено в оценке: метод не обещает чуда на вырожденном окне. На практике глубину окна берут небольшой (\(p\) порядка единиц), и обусловленность контролируема; при необходимости старые направления выбрасываются. Это же объясняет умеренную деградацию L-PB по глубине памяти на слайде 11.\par
\end{spoken}

\textbf{В: Почему коррекция SS-U именно ортогональна \(s_{k}\)?}
\begin{spoken}
Потому что это \textbf{единственная} степень свободы, которая не разрушает уже достигнутое текущее секущее условие. Любая компонента вдоль \(s_{k}\) сместила бы \(B_{k+1} s_{k}\) с \(y_{k}\). Ортогональное дополнение — это ровно то место, где можно «бесплатно» согласовываться с прошлыми секущими, сохраняя симметрию.\par
\end{spoken}

\textbf{В: SS-SR1 не даёт выигрыша — зачем тогда конструкция?}
\begin{spoken}
Затем, что SR1 со skip-rule \textbf{сам по себе не деградирует} — ему нечего стабилизировать. Конструкция ценна там, где базовое обновление деградирует, — как PSB. Я показываю и то, и другое, чтобы честно очертить, \textbf{когда} метод полезен, а не маскировать отсутствие эффекта.\par
\end{spoken}

\textbf{В: Какой ваш личный вклад в главу 4, если статья соавторская?}
\begin{spoken}
Сама оценка скорости метода VIJI и нижняя оценка — результат соавторской работы (NeurIPS 2024), я этого не присваиваю. \textbf{Мой вклад в диссертации} — Утверждение 4.1: подстановка моей локальной мультисекущей оценки в параметр неточности \(\delta \), что заменяет грубый равномерный потолок на подпространственный \(O(\rho _{k})\)-контроль. То есть глава 4 — это мост между моей локальной теорией и внешним методом.\par
\end{spoken}

\textbf{В: Сравнивали ли с L-BFGS по времени, а не по итерациям?}
\begin{spoken}
По итерациям и устойчивости — да, поведение L-PB согласуется с L-BFGS. Сравнение по wall-clock на индустриальных задачах — честно, это в перспективах; текущая реализация исследовательская, и основной вклад работы аналитический.\par
\end{spoken}

\textbf{В: Воспроизводимость?}
\begin{spoken}
Вся реализация и скрипты построения всех рисунков выложены в открытый доступ — это отмечено и в отзыве руководителя.\par
\end{spoken}

\medskip\noindent{\color{rulecol}\rule{\linewidth}{0.4pt}}\par\medskip

\section*{\color{headcol} ЧАСТЬ 5. ЗАКЛЮЧИТЕЛЬНОЕ СЛОВО (если попросят)}

\begin{spoken}
Благодарю комиссию за вопросы, научного руководителя Дмитрия Игоревича Камзолова — за постановку задачи и поддержку, и рецензента Александра Викторовича Рогозина — за внимательный разбор. Замечания принимаю и вижу в них конкретный план дальнейшей работы. Спасибо.\par
\end{spoken}

\end{document}